\section{Experimental protocol}
\label{sec:protocol}
\paragraph{Data.}
We use the binary item-level matrices released with~\citet{zhang2025misses} (fixed commit \texttt{55faff4}, MIT licence), derived from the Open LLM Leaderboard v2~\citep{fourrier2024openllm} and HELM~\citep{liang2023helm}: seven Open LLM Leaderboard benchmarks (ARC-Challenge, BBH, GPQA, IFEval, MMLU-Pro, MuSR, MATH Lvl~5; 218--448 models, 541--12\,032 items) and six HELM scenarios (files named commonsense, gsm, legalbench, math, med\_qa, mmlu; 79--83 models, 437--2\,047 items). Model sets differ per benchmark (Appendix Table~\ref{tab:data}). We did not re-run any model.

\paragraph{Splits (who is the target?).}
For each benchmark and seed we split models into source and target sets in one of three ways. \textbf{Random}: 30\% of models are targets, drawn uniformly (the interpolation setting most prior work evaluates). \textbf{Held-out organisations}: whole organisations (the prefix of the model id; Open LLM Leaderboard only, since HELM model ids are not aligned with all matrices) are moved to the target set until $\ge25\%$ of the models are targets; this removes near-duplicate fine-tunes and merges of the same base model from the sources. \textbf{Stronger targets}: the 25\% of models with the highest full-benchmark accuracy are targets, all others are sources. This is a deliberately demanding but realistic extrapolation (a new, better model is evaluated using the records of older, weaker ones); the split itself is deterministic and seeds vary the coreset and model initialisation. Extra severities (top 10\%, top 50\%, weakest 25\%) are in Appendix~\ref{app:shift}.

\paragraph{Budgets, seeds, metrics.}
Coreset sizes $n\in\{10,25,50,100,200\}$; 30 seeds per (benchmark, split type) (random splits differ per seed); all methods share the same split within a seed. For each run we compute, over the target models of the split, RMSE, MAE and signed mean error of $\hat a_t-a_t$ (accuracy in $[0,1]$), and Spearman's $\rho$ and Kendall's $\tau$ between $\hat a_t$ and $a_t$ (note that with a shared coreset the common component of the error cancels in rankings). To compare across benchmarks with very different scales (e.g.\ MATH Lvl~5 accuracies are around 5\%) we report \emph{relative RMSE}: for a benchmark, split type and $n$, $\rho_b=\sqrt{\overline{\mathrm{MSE}}_{\text{method}}/\overline{\mathrm{MSE}}_{\text{random sampling}}}$ with MSE averaged over seeds, aggregated across benchmarks by the geometric mean (below 1 = better than random sampling with the same $n$). 95\% confidence intervals come from a hierarchical bootstrap that resamples benchmarks and then seeds within benchmarks (1000 resamples), jointly for all methods (so differences between methods are paired). Intervals reflect variability over benchmarks, splits and coresets; with the deterministic stronger-targets split they do \emph{not} reflect variability in which models happen to be the strongest.

\paragraph{Baselines and hyper-parameters.}
Baselines are random sampling, difficulty-stratified sampling, linear calibration, ridge regression, Anchor Points, anchors with IRT plug-in (tinyBenchmarks-style), D-optimal items with IRT plug-in (information-based), and IRT-AIPW with uniform sampling (the AIPW-style estimator that~\citet{zhang2025misses} recommend; our implementation, not theirs). The IRT dimension is fixed to $d=4$ in the main tables ($d\in\{2,4,8\}$ in ablations). The IRT regularisers, the ability prior variance (4), the Platt penalty (1) and $c_0=0.5$ were fixed using a separate development run (seeds $\ge1000$ and 2000--2007, on 9--11 benchmarks, budgets 25 and 100; the development benchmarks did not include BBH or MMLU-Pro) and never tuned on the evaluation seeds. The main variants (VA-PPS $\lambda=1$ with Cal-IRT and with ShrinkCal-IRT, $d=4$) were declared before the main run; the many other combinations in the ablations are reported for completeness and not selected post hoc.
